\infographicchapter{Ultrasound Elastography}{ultrasound_elastography.png}
\label{ch:infographic-ultrasound-elastography}

\section{Overview}

Ultrasound elastography is a gentle, painless ultrasound test that measures how soft or stiff tissues are inside your body---working like a high-tech, deeply accurate version of a doctor feeling for a firm lump by hand. Because scarring (fibrosis) and many tumors make tissue firmer than healthy surrounding tissue, measuring stiffness gives doctors important clues without using needles. Two main methods are used: strain elastography, which watches how much tissue squishes under gentle probe pressure, and shear-wave elastography (SWE), which sends a tiny, harmless vibration or focused sound ripple sideways through the tissue and times how fast it travels---faster ripples mean stiffer tissue, reported in meters per second or kilopascals (kPa).

The test layers a color-coded stiffness map right on top of a regular black-and-white ultrasound image. Its most common and proven use is checking for liver scarring (fibrosis) and cirrhosis in chronic liver conditions such as viral hepatitis or fatty liver disease---often sparing patients from an invasive liver biopsy. Doctors also use elastography alongside regular ultrasound to check the firmness of breast lumps, thyroid nodules, tendons and muscles, and occasionally the prostate or kidneys.

Ultrasound elastography is fast, widely available, completely free of X-ray radiation, and fits smoothly into a routine ultrasound visit. Its main limits are that results can vary with examiner technique (especially in strain mode), thick tissue layers in obesity can weaken the signal, and different scanner brands use slightly different stiffness cutoff numbers. Also, temporary liver inflammation, fluid buildup, or a recent meal can make the liver look stiffer than it really is, so doctors always read stiffness numbers alongside regular ultrasound and lab results.

\section{Underlying Physical Principles}

\begin{itemize}
    \item \textbf{Creating a Tiny Tissue Nudge:} Elastography gently nudges tissue in one of three ways: light hand pressure with the probe (strain elastography), a harmless focused Acoustic Radiation Force Impulse (ARFI) push pulse---measuring either a single targeted spot (point shear-wave elastography, pSWE) or a wider real-time box (2D shear-wave elastography, 2D-SWE)---or a soft mechanical tap on the side of the ribs (transient elastography, such as FibroScan).
    \item \textbf{Faster Ripples Mean Stiffer Tissue:} That tiny nudge sends microscopic sideways ripples (called shear waves) spreading through the organ, much like ripples on a pond. Ripples travel slowly through soft, healthy tissue and much faster through scarred or firm tissue. The scanner converts ripple speed into an absolute stiffness score called Young's modulus in kilopascals (kPa), extending manual palpation deep inside the body.
    \item \textbf{Ultra-Fast Echo Tracking:} Immediately after the nudge, the probe fires rapid tracking sound pulses to follow the tiny ripples moving across the tissue, measuring movements smaller than the width of a human hair.
    \item \textbf{Drawing a Color Stiffness Map:} The computer calculates local ripple speed across the viewing box and paints a smooth, color-coded stiffness map right over your standard black-and-white ultrasound picture.
\end{itemize}

\section{Imaging Mechanism}

\begin{enumerate}
    \item \textbf{Applying Gel and Choosing a Clean Spot:} The sonographer applies generous ultrasound gel, locates the organ or lump on regular black-and-white ultrasound, and places a small sampling box away from large blood vessels and organ edges.
    \item \textbf{Sending the Gentle Push or Tap:} The scanner delivers a brief, painless acoustic push pulse, a soft 50 Hz mechanical tap (in transient mode), or gentle hand pressure to launch microscopic sideways ripples through the tissue.
    \item \textbf{Tracking the Ripples at High Speed:} The probe immediately switches to ultra-fast listening mode to capture returning echoes thousands of times per second as the ripples travel sideways.
    \item \textbf{Calculating Stiffness Scores:} Built-in software tracks the ripple speed across the box, filters out noisy readings, and calculates the median stiffness in meters per second or kilopascals (kPa) along with a quality-check score.
    \item \textbf{Displaying the Color Map and Report:} A translucent color stiffness map appears directly over the grayscale anatomy, alongside clear numerical summaries for the radiologist or liver specialist.
\end{enumerate}

\section{Key Technical Parameters}

\noindent
\begin{tabularx}{\textwidth}{@{}l X@{}}
\toprule
\textbf{Parameter} & \textbf{Typical Value / Description} \\
\midrule
Physical Principle & Measuring tissue squishiness (strain) or sideways ripple (shear-wave) speed in meters per second or kilopascals (kPa) \\
Energy / Signal Source & Focused ultrasound push pulses, gentle mechanical tap, or light probe pressure tracked by diagnostic ultrasound \\
Typical Frequency / Energy & 1--5 MHz (abdominal probes) or 50 Hz soft tap (transient liver scan); 6--15 MHz (breast, thyroid, and tendon probes) \\
Spatial Resolution & About 1--3 mm on 2D color stiffness maps; 10--40 mm sampling zone for point or transient liver checks \\
Typical Scan Duration & 5--15 minutes for elastography alone (15--30 minutes with a full ultrasound) \\
Contrast Agent & None needed \\
Ionizing Radiation & None (0 mSv) \\
\bottomrule
\end{tabularx}

\section{Common Symptoms That Prompt This Scan}
\begin{itemize}
    \item \textbf{Often Silent---Abnormal Liver Tests or Fatty Liver:} Liver scarring rarely causes early symptoms, so liver elastography is most often prompted by fatty liver, viral hepatitis, or elevated liver enzymes.
    \item \textbf{A New Breast or Thyroid Lump:} Evaluating whether a lump seen on standard ultrasound is soft (more likely benign) or stiff (more likely to need a biopsy).
    \item \textbf{Persistent Fatigue or Upper-Right Belly Discomfort:} Non-specific symptoms in people with chronic liver risk factors needing a quick, painless stiffness check instead of a liver biopsy.
    \item \textbf{Tracking Chronic Liver or Muscle Conditions:} Monitoring whether liver stiffness is improving with treatment, or assessing stubborn tendon and muscle stiffness.
    \item \textbf{Liver Risk Factors and Late Warning Signs:} Proactive screening for people with type 2 diabetes, autoimmune or genetic liver disorders, or long-term use of liver-stressing medicines (such as methotrexate), as well as evaluating cirrhosis signs like jaundice (yellowing of the skin or eyes), easy bruising, dark urine, or swelling in the legs or belly.
\end{itemize}

\section{Clinical Usages}

\subsection{Primary Usages}

\begin{itemize}
    \item \textbf{Checking Liver Scarring (Fibrosis and Cirrhosis):} Measuring liver stiffness noninvasively to check for early scarring or advanced cirrhosis in viral hepatitis, metabolic fatty liver disease (MASLD), and alcohol-related liver disease.
    \item \textbf{Checking Portal Vein Pressure Risk:} Measuring both liver and spleen stiffness to estimate the risk of high pressure in the liver's main vein (portal hypertension) and swollen food-pipe veins (varices), helping decide who needs an endoscopy.
    \item \textbf{Evaluating Breast Lumps:} Adding stiffness clues to regular breast ultrasound to help tell firmer cancers apart from softer harmless fibroadenomas or cysts, reducing unnecessary biopsies.
    \item \textbf{Evaluating Thyroid Nodules:} Helping spot firmer thyroid nodules that are more likely to be cancer so doctors know which ones truly need a fine-needle biopsy.
    \item \textbf{Checking Tendons and Muscles:} Measuring elasticity in Achilles or knee tendons, the bottom of the foot (plantar fascia), and shoulder rotator cuff tendons to assess wear or healing.
\end{itemize}

\subsection{Secondary Usages}

\begin{itemize}
    \item \textbf{Tracking Liver Treatment Progress:} Repeating liver stiffness checks over the years to see whether scarring is staying stable or improving after hepatitis treatment, weight loss, or new liver medicines.
    \item \textbf{Prostate and Testicular Checks:} Helping spot firmer areas inside the prostate or testicle during ultrasound to guide targeted biopsies.
    \item \textbf{Kidney Transplants and Chronic Kidney Disease:} Checking for early kidney tissue scarring alongside standard Doppler ultrasound.
    \item \textbf{Lymph Nodes and Salivary Glands:} Helping tell harmless swollen lymph nodes apart from cancerous ones, and checking gland stiffness in chronic salivary gland conditions such as Sjogren syndrome.
\end{itemize}

\section{Advantages}

\begin{itemize}
    \item \textbf{Needle-Free and Radiation-Free:} Measures tissue firmness with zero needles, zero X-rays, and zero contrast dye, dramatically cutting down the need for invasive liver biopsies.
    \item \textbf{Built Right Into Routine Ultrasound:} Takes just a few extra minutes on modern ultrasound machines, letting doctors view stiffness maps right on top of regular anatomy pictures.
    \item \textbf{Accessible and Affordable:} Widely available in clinics at much lower cost and shorter visit time than MRI elastography (MRE).
    \item \textbf{Objective Numbers for Tracking Over Time:} Shear-wave modes provide clear numerical readings (m/s or kPa) that help track liver health or treatment response year by year.
    \item \textbf{Sharper Lump Assessment:} Adds valuable firmness information to standard ultrasound when evaluating breast, thyroid, and surface soft-tissue lumps.
\end{itemize}

\section{Disadvantages}

\begin{itemize}
    \item \textbf{Sensitive to Technique:} Strain elastography relies on steady hand pressure, and even shear-wave readings can shift if the probe presses too hard, sits at an angle, or you take a deep breath during the reading.
    \item \textbf{Harder in Deep Tissue or Belly Fluid:} Severe obesity, thick skin layers, narrow rib spaces, or fluid around the liver (ascites) can block or weaken the ripples and prevent reliable readings.
    \item \textbf{Temporary Conditions Can Skew Liver Numbers:} Active liver inflammation (hepatitis), blocked bile ducts, heart-related liver congestion, or eating a meal shortly beforehand can temporarily raise stiffness and mimic cirrhosis.
    \item \textbf{Numbers Vary Across Scanner Brands:} Different ultrasound makers and modes use slightly different cutoff numbers, so repeat liver scans work best on the same brand of machine.
    \item \textbf{Stiffness Alone Is Not a Final Diagnosis:} Some harmless calcified or fibrous lumps feel stiff, while certain soft cancers feel supple---so elastography supports, rather than replaces, regular imaging or biopsy.
\end{itemize}

\section{Contraindications}

\subsection{Absolute}

\begin{itemize}
    \item \textbf{None:} Because ultrasound elastography is completely noninvasive, uses no dye, and uses no radiation, there are no strict patient contraindications.
\end{itemize}

\subsection{Relative}

\begin{itemize}
    \item \textbf{Conditions That Skew Liver Readings:} Large amounts of belly fluid around the liver (which blocks transient elastography ripples), severe active hepatitis flares, blocked bile ducts, or uncontrolled right-sided heart failure can make liver stiffness numbers unreliable.
    \item \textbf{Delicate Areas and Skin Limits:} Examiners avoid aiming repeated high-energy push pulses at a developing baby, the eyes, or air-filled lungs; open wounds, painful broken ribs, or skin infections right over the scan spot may mean postponing the test.
\end{itemize}

\section{Patient Safety Guidelines}

\begin{itemize}
    \item \textbf{Keeping Sound Energy Safe (ALARA):} Because shear-wave push pulses use brief bursts of slightly higher energy than standard ultrasound, sonographers follow AIUM and WFUMB guidelines, monitor on-screen Thermal Index (TI) and Mechanical Index (MI) gauges, and limit repeat pulses.
    \item \textbf{Protecting Sensitive Organs:} High-energy push pulses are kept away from a developing fetus, eyes, and air-filled lungs unless specifically approved under strict safety limits.
    \item \textbf{Fasting and Rest Prevent False Alarms:} Fasting for 4--6 hours and resting quietly beforehand keeps normal mealtime blood flow from temporarily stiffening the liver and triggering unnecessary worry or biopsies.
    \item \textbf{Light Touch and Clean Probes:} Sonographers use generous gel and very light probe contact to keep you comfortable and prevent artificially squeezing shallow lumps, and disinfect probes thoroughly between every patient.
    \item \textbf{Consistent Tracking:} Reports record the exact scanner brand and mode used so your doctor can compare future readings accurately and safely.
\end{itemize}

\section{Prerequisite Tests}

\begin{itemize}
    \item \textbf{No Blood Tests Needed for Safety (Recent Liver Labs Helpful):} No lab work is required to make the scan safe, though your doctor will review recent liver blood tests (AST, ALT, bilirubin) before liver elastography because active inflammation or blocked bile can falsely raise stiffness.
    \item \textbf{Standard Ultrasound First:} A regular black-and-white and Doppler ultrasound is done first to pick a clean spot away from big blood vessels and check for belly fluid or liver congestion.
    \item \textbf{Fasting Check and History:} Staff confirm you fasted 4--6 hours before a liver scan and review your liver history or prior breast and thyroid scans.
    \item \textbf{Checking Past Scanner Brand:} If you are tracking liver scarring over time, staff check which scanner brand and mode were used previously so results line up cleanly.
\end{itemize}

\section{Preparing for the Scan}

\begin{itemize}
    \item \textbf{Fast Before Liver Elastography:} Do not eat food or drink anything with calories for 4--6 hours before a liver stiffness scan, because digesting a meal boosts blood flow to the liver and artificially raises stiffness. Small sips of plain water and essential pills are fine.
    \item \textbf{Rest Quietly Beforehand:} Avoid heavy exercise right before your appointment and sit calmly for 10--15 minutes before a liver scan so blood flow settles to a steady resting level.
    \item \textbf{Easy-Access Clothing:} Wear loose, comfortable clothes so the right side of your ribcage (for liver scans) or your breast, neck, or joint can be uncovered easily, and skip thick skin lotions or oils over the area.
    \item \textbf{Easy Position and Gentle Breath Pause:} For liver scans, you will lie on your back with your right arm raised behind your head to open the spaces between your ribs, and practice pausing your breath gently midway---without taking a huge deep breath or straining.
    \item \textbf{Stay Relaxed for Shallow Scans:} For breast, thyroid, or tendon scans, simply relax your muscles and stay still while the examiner uses plenty of gel and a feather-light touch.
\end{itemize}

\section{During Scan}

\begin{itemize}
    \item \textbf{Comfortable Positioning:} For a liver scan, you lie on your back with your right arm tucked comfortably behind your head to widen the rib spaces. For breast, thyroid, or joint scans, you lie or sit with the area relaxed and supported.
    \item \textbf{Gentle Breath Pauses:} For liver readings, the sonographer coaches you to pause breathing comfortably mid-breath for 3--5 seconds per measurement---avoiding deep breaths or bearing down, which can skew stiffness numbers.
    \item \textbf{What It Feels Like:} You feel cool gel and very light, steady probe contact. During transient liver elastography (like FibroScan), you feel a gentle, painless flick or tap against your right ribs; during shear-wave elastography, you cannot feel the sound ripples at all.
    \item \textbf{Speak Up Anytime:} Tell the sonographer if raising your arm, pausing your breath, or probe contact over a tender rib, breast, or joint feels uncomfortable so they can adjust right away.
    \item \textbf{Quick and Easy Visit:} Stiffness measurements usually take just 5--15 minutes (or 15--30 minutes if combined with a full regular ultrasound), with several quick repeat readings to ensure accuracy.
\end{itemize}

\section{Normal Appearance}

\begin{itemize}
    \item \textbf{Soft, Healthy Liver:} Shows an even, low-stiffness color map (commonly smooth blue) with median stiffness below about 5 kPa (ripple speed under 1.3 m/s) and tight consistency across repeat readings.
    \item \textbf{Normal Breast and Thyroid Tissue:} Healthy breast tissue, fat, and thyroid lobes look evenly soft-to-medium in color without hard focal spots or stiff outer rings.
    \item \textbf{Healthy Tendons and Muscles:} Healthy tendons are naturally taut and firm along their length (fast ripple speed), while relaxed muscles look softer and even.
\end{itemize}

\section{Observing for Abnormal Patterns}

\begin{itemize}
    \item \textbf{Stiffer Liver Indicating Scarring:} Higher liver stiffness reflects increasing scar tissue (fibrosis). As a general guide, readings below about 5 kPa are normal, while readings above 13--15 kPa strongly point to advanced liver scarring or cirrhosis. Reports often summarize this on a scale from F0--F1 (little or no scarring) through F2 (moderate) and F3 (severe) to F4 (cirrhosis), alongside a CAP score for liver fat (5\% or less is normal).
    \item \textbf{Firm Breast or Thyroid Lumps:} Cancers in the breast or thyroid often appear noticeably stiff on color maps (frequently above 70--80 kPa in breast masses) and may show a stiff rim extending slightly beyond the lump's black-and-white border. On strain elastography, doctors also look at the strain ratio (comparing how little a lump squishes next to nearby soft fat) or a 1-to-5 color stiffness scale such as the Tsukuba score.
    \item \textbf{Classic Soft and Cystic Signs:} Simple fluid-filled cysts show a telltale layered ``bull's-eye'' pattern on strain elastography or a dark empty center on shear-wave scans (since sideways ripples do not travel through simple liquid), while harmless fibroadenomas and fatty lipomas usually look evenly soft.
    \item \textbf{Softer Tendons in Wear or Tears:} Unlike scarred livers or tumors, worn or partially torn tendons actually become \emph{softer} than healthy tight tendon fibers, showing slower ripple speeds.
    \item \textbf{Common False-High Mimics:} Active hepatitis flares, blocked bile ducts, heart-related liver congestion, or eating shortly before the scan can make a liver look falsely cirrhotic; pressing too hard with the probe or scanning a hard calcium speck can make a harmless breast or thyroid lump look falsely stiff.
\end{itemize}

\section{Dealing with Uncertainties}

\begin{itemize}
    \item \textbf{Temporary Liver Stiffness Boosters:} Active liver inflammation, blocked bile, heart congestion, or a recent meal can push stiffness numbers up and mimic permanent scarring.
    \item \textbf{Technical Limits:} Severe obesity, narrow rib spaces, fluid around the liver, or deep targets can weaken ripples and cause scattered readings (high variability above 30\%).
    \item \textbf{Brand Differences and Overlapping Lumps:} Stiffness cutoffs differ slightly across scanner brands, and harmless calcified lumps can feel stiff while rare soft cancers can feel supple.
    \item \textbf{Quick Repeat Check:} If readings bounce around or probe pressure seemed too firm, the sonographer can confirm fasting, relax your breath-hold, pick a cleaner window away from big vessels, and repeat a few readings.
    \item \textbf{Confirming With Other Tests:} If liver numbers disagree with your symptoms or blood work, your doctor may recheck after inflammation cools down, combine results with blood scarring scores (like FIB-4), or order an MRI elastography (MRE). Uncertain breast or thyroid lumps follow standard BI-RADS or TI-RADS rules and may be sampled with a quick needle biopsy.
\end{itemize}

\section{Augmenting Methods}

\begin{itemize}
    \item \textbf{MRI Elastography (MRE):} The most accurate noninvasive backup for liver scarring when severe obesity, narrow ribs, belly fluid, or unclear ultrasound numbers limit the view, sampling a much larger portion of the liver.
    \item \textbf{Full Ultrasound and Fat Measurement:} Combining stiffness checks with regular ultrasound, Doppler blood flow, and ultrasound fat-measuring tools (CAP or attenuation imaging) checks liver fat, blood flow, and lumps all in one visit.
    \item \textbf{Blood Liver Panels and Scarring Scores:} Pairing stiffness numbers with AST, ALT, bilirubin, platelets, and simple blood scores (FIB-4, APRI, or ELF) rules out temporary inflammation and sharpens cirrhosis risk checks.
    \item \textbf{Mammograms, Breast MRI, and Needle Biopsy:} Combining elastography with mammography, MRI, and BI-RADS or TI-RADS scoring guides accurate biopsies when needed.
    \item \textbf{Upper Endoscopy and Liver Biopsy:} When high liver and spleen stiffness point to advanced liver disease, an upper endoscopy checks for swollen food-pipe veins, while a liver biopsy remains the final arbiter when the cause of liver disease is still unclear.
\end{itemize}

\section{Comparison with Alternative Modalities}

\begin{itemize}
    \item \textbf{Versus MRI Elastography (MRE):} MRE samples a wider liver area and works better in severe obesity or belly fluid, while ultrasound elastography is much faster, far less expensive, and available right in clinics---making ultrasound the first-line choice and MRE the backup for tough cases.
    \item \textbf{Versus Liver Biopsy:} A needle biopsy gives microscopic cell detail when diagnosis is unclear, but carries small bleeding risks and checks only a tiny sliver of liver. Ultrasound elastography is completely painless and ideal for routine staging and yearly tracking.
    \item \textbf{Versus Regular Ultrasound Alone:} Standard ultrasound spots advanced cirrhosis or obvious lumps but misses early-to-moderate liver scarring; adding elastography measures stiffness during the very same visit.
\end{itemize}

\section{Follow-up Steps}

\begin{itemize}
    \item \textbf{Right After the Scan:} There is zero recovery time---you can eat, drink, and return to normal activities immediately. Before you leave, the sonographer confirms the readings met strict quality checks.
    \item \textbf{Reviewing With Your Blood Tests:} Your doctor looks at stiffness numbers together with your regular ultrasound and liver blood tests. If temporary inflammation, blocked bile, or a recent meal skewed the result, the scan is simply repeated once things settle.
    \item \textbf{Next Steps If Stiffness Is High:} High liver stiffness may lead to blood scarring scores, an endoscopy check for varices, or MRE; stiff breast or thyroid lumps are evaluated using BI-RADS or TI-RADS guidelines to decide whether a biopsy is needed.
    \item \textbf{Tracking on the Same Platform:} For chronic liver conditions, follow-up scans are commonly repeated every 1--3 years---ideally on the same brand and type of elastography machine for clean year-to-year comparisons.
    \item \textbf{Treating the Cause Helps Protect Your Liver:} When liver fibrosis or excess fat is caught early, treating the underlying condition and adopting healthy diet and exercise habits can stop scarring from worsening and may even allow some scar tissue to slowly improve.
\end{itemize}

\section{Key References}

\begin{itemize}
    \item Barr RG, Wilson SR, Rubens D, Garcia-Tsao G, Ferraioli G. Update to the Society of Radiologists in Ultrasound Liver Elastography Consensus Statement. \textit{Radiology}. 2020;296(2):263--274.
    \item Ferraioli G, Barr RG, Farrokh A, et al. How to Perform Shear Wave Elastography. Part I and Part II. \textit{Medical Ultrasonography}. 2022;24(1):95--106.
    \item Barr RG, Nakashima K, Amy D, et al. WFUMB Guidelines and Recommendations for Clinical Use of Ultrasound Elastography: Part 2: Breast. \textit{Ultrasound in Medicine and Biology}. 2015;41(5):1148--1160.
    \item Sigrist RMS, Liau J, Kaffas AE, Chammas MC, Willmann JK. Ultrasound Elastography: Review of Techniques and Clinical Applications. \textit{Theranostics}. 2017;7(5):1303--1329.
    \item National Library of Medicine (MedlinePlus). \textit{Elastography: MedlinePlus Medical Test}. National Institutes of Health. Available at: \url{https://medlineplus.gov/lab-tests/elastography/}.
    \item Radiopaedia.org. \textit{Elastography}. Radiopaedia. Available at: \url{https://radiopaedia.org/articles/elastography}.
\end{itemize}

\section{Further Reading}
\begin{itemize}
  \item Wikipedia contributors. \textit{Elastography}. Wikipedia, The Free Encyclopedia. Available at: \url{https://en.wikipedia.org/wiki/Elastography}.
\end{itemize}
